\chapter{Introduction}
\label{chap:introduction}

\section{Motivation: Decoding Affect from Neurophysiology as an Electrical Engineering Problem}

This dissertation is an electrical engineering study of structured, noisy, nonstationary biosignals. The application domain is affective neurophysiology, but the technical problem is signal representation under measurement noise, physiological artifacts, between-participant variability, and uncertain spatial mixing. Scalp electroencephalography (EEG) is therefore treated as a multichannel observation of distributed cortical activity, not as a direct readout of mental state. The study asks whether a fixed nonlinear transformation can produce useful and inspectable summaries of these observations. Evaluation includes prediction, sensitivity analyses, and examination of intermediate quantities. Clinical actionability is outside the scope of the study.

The governing question is whether a nonlinear transformation can expose stable, compressible, and structurally coherent patterns in noisy neural measurements. The proposed staged architecture combines a spiking reservoir, a temporal coding bottleneck, and graph-organized spatial analysis. Classification accuracy is one evaluation criterion; the others are operating-regime stability, compression behavior, representation geometry, dynamical response properties, and cross-scale correspondence. The analysis asks which intermediate quantities remain traceable to defined parts of the input and where that traceability breaks down.

\subsection{The Engineering Difficulty of Scalp EEG}

Scalp EEG records microvolt-scale voltage fluctuations produced by spatially mixed cortical sources. Volume conduction through cerebrospinal fluid, skull, and scalp makes the source-to-sensor mapping many-to-one and ill-conditioned \cite{nunez_srinivasan_2006}. EEG offers millisecond temporal sampling but limited source localization without an explicit forward model. It is therefore misleading to rank its ``spatial resolution'' against fMRI with a single orders-of-magnitude statement; the modalities measure different physical quantities and depend on different acquisition and inverse-model choices.

Ocular and muscular artifacts, cardiac interference, line noise, and electrode-impedance drift can be large relative to event-related potentials (ERPs) of interest \cite{nunez_srinivasan_2006}. ERP amplitude and latency also vary across trials and participants. Differences in anatomy and electrode placement add further between-participant variability.

EEG is also nonstationary: its statistical properties vary over time with behavioral state, task engagement, and measurement conditions. Event-related responses are transient and temporally localized. These properties motivate explicit temporal processing and evaluation across participants rather than an assumption that every epoch is an independent draw from one fixed distribution.

\subsection{Clinical Context and Translation Limits}

ERP studies report group-level associations between selected components and psychiatric symptoms or diagnoses \cite{gehring_erp_2012,nelson_hajcak_2017}. Such associations do not automatically yield individual-level diagnostic or prognostic tools. Reliability, heterogeneity, confounding, multiplicity, and external validation all matter, particularly when analyses use trial-averaged rather than single-trial responses \cite{weinberg_hajcak_2011}.

The SHAPE cohort also has substantial comorbidity and medication use. Diagnostic labels are therefore treated as overlapping metadata, not as mutually exclusive biological states. Clinical-label analyses in this dissertation are exploratory screens; they do not establish biomarkers, mechanisms, or individual-level clinical utility.

\subsection{The SHAPE Study as the Principal Validation Setting}

The real-data application uses the Stress, Health, and the Psychophysiology of Emotion (SHAPE) study \cite{noauthor_shape_2018}. The analyzed sample contains $211$ participants with overlapping diagnostic and medication metadata. Participants viewed International Affective Picture System images \cite{lang_iaps_2008}; the data supplied for this work are participant-by-condition trial-averaged ERP matrices for Negative, Neutral, and Pleasant conditions.

The repeated three-condition design supports within-participant comparisons, and the $34$ channels permit channel-wise feature analysis. Important limits are equally central: trial averaging removes single-trial variability; electrode names and montage coordinates were unavailable in the supplied analysis files; clinical groups overlap; and the raw SHAPE EEG is not included with the dissertation source package. The post-defense audit therefore could not recompute Chapter~5's corrected fold-local PCA results.

\subsection{Why EEG? An Engineering Justification}

The motivation for using EEG is engineering as much as physiological. EEG provides millisecond temporal sampling and a natural multichannel structure while presenting low signal-to-noise ratio, nonstationarity, between-participant heterogeneity, and uncertain spatial mixing. These properties motivate evaluation of temporal coding, dimensionality reduction, robustness, and relationships among channel features.

The engineering scope of scalp EEG is well defined. It provides millisecond temporal resolution at the cost of spatial precision relative to fMRI, source-level specificity relative to intracranial recordings, and metabolic information relative to positron emission tomography. The dissertation's signal-processing contributions operate within the information content of the scalp-recorded signal, and downstream clinical interpretation respects these measurement-model boundaries.

\subsection{Formal Problem Statement: Affective EEG as an Inverse Problem}
\label{sec:inverse_problem}

The preceding discussion motivates a precise mathematical formulation of the inference problem that ARSPI-Net is designed to address. This formulation unifies the engineering, clinical, and theoretical threads of the dissertation into a single inverse-problem framework.

Let $\mathcal{P}$ denote the population of individuals, each indexed by $p \in \mathcal{P}$. For every individual~$p$:
\begin{itemize}
    \item Let $D_p$ denote unobserved neural-source dynamics, without assuming that they can be identified from scalp EEG.
    \item Let $\delta_p\in\Delta$ denote observed, overlapping clinical labels. These labels are associated metadata, not assumed causal parameters of $D_p$.
\end{itemize}

A known affective stimulus $s \in S$ (an IAPS image assigned to an experimental condition) is presented. Although IAPS includes normative ratings, stimulus- and participant-level valence or arousal scores are not modeled in this dissertation. The internal neural response is represented abstractly as
\begin{equation}
    r_p(s) = \mathcal{F}(D_p, s),
    \label{eq:forward_dynamics}
\end{equation}
where $\mathcal{F}$ denotes an unknown stimulus-response process. The dissertation does not estimate $\mathcal{F}$ or infer diagnosis from this equation.

The experimenter observes only the scalp EEG/ERP signal
\begin{equation}
    x_p = g(r_p) + \eta,
    \label{eq:observation}
\end{equation}
where $g$ is the forward volume-conduction operator (mapping cortical sources to scalp potentials) and $\eta$ encompasses measurement noise and physiological artifacts. In the SHAPE dataset, trial averaging and standard preprocessing have already been applied, so the observed $x_p(s) \in \mathbb{R}^{T \times C}$ ($T$ time samples, $C = 34$ channels) represents the scalp-level ERP matrix for individual~$p$ under stimulus condition~$s$. The volume-conduction operator $g$ is embedded in the observation; the dissertation works with the scalp-level signal directly rather than attempting cortical source reconstruction.

Because $g$ is not inverted and source identifiability is not established, the dissertation does not claim to recover $D_p$. Its operational question is:
\begin{equation}
    \text{Construct and evaluate model summaries of } x_p(s) \text{ across conditions and participants.}
    \label{eq:inverse_problem}
\end{equation}
Equivalently, the goal is to construct a mapping
\begin{equation}
    \phi: x_p(s) \;\mapsto\; \mathbf{z}_p(s),
    \label{eq:phi}
\end{equation}
where $\mathbf{z}_p(s)$ contains reservoir, embedding, and graph-derived summaries. The mapping is evaluated against three requirements:
\begin{enumerate}
    \item \textbf{Computational stability}: sensitivity to selected reservoir parameters, initializations, and synthetic perturbations is measured.
    \item \textbf{Traceability}: each implemented transformation is specified, and any relationship to ERP quantities is treated as an empirical association rather than assumed meaning.
    \item \textbf{Cohort-level evaluation}: condition classification and clinical-label associations are evaluated with participant grouping and explicit inferential limitations.
\end{enumerate}

ARSPI-Net instantiates $\phi$ as a staged pipeline: a fixed spiking reservoir transforms $x_p(s)$; BSC$_6$ summarizes early spike activity; PCA compresses features; and subsequent analyses compute classifier, dynamical, and graph summaries. Condition labels are used for prediction and repeated-measures comparisons. Clinical labels are used only in exploratory cohort analyses.

Each subsequent chapter addresses one part of this construction. Chapters~\ref{chap:lsm} and~\ref{chap:embeddings} use controlled synthetic experiments to select and characterize the reservoir and code. Chapter~\ref{chap:arspi_net} reports the SHAPE application and the post-defense preprocessing limitations. Chapter~\ref{chap:dynamics} defines reservoir-response summaries, and Chapter~\ref{chap:synthesis} reports descriptive associations among model-derived quantities. Chapter~\ref{chap:conclusion} consolidates the supported findings and unresolved tests.


\section{The Representation Gap and the Neuromorphic Opportunity}

\subsection{Motivating an Alternative to Conventional Deep Learning for Biosignals}

Architectures such as LSTMs and Transformers have demonstrated substantial success in generic time-series modeling \cite{ismail_fawaz_deep_2019}. When applied to physiological signals, however, several structural properties of these architectures motivate the investigation of alternative paradigms.

First, dense neural models can require substantial memory traffic and multiply--accumulate operations. Event-driven hardware offers a different engineering trade-off for some workloads. This dissertation does not benchmark hardware, edge deployment, or energy consumption, so efficiency is a motivation rather than a demonstrated advantage.

Second, EEG nonstationarity and between-participant variation complicate model evaluation. Third, different architectures expose different intermediate quantities. These observations motivate comparing a fixed spiking reservoir with conventional baselines; they do not imply that LSTMs, Transformers, or other models are intrinsically unsuitable for EEG.

These observations motivate an investigation of one alternative representation: a fixed spiking reservoir followed by conventional statistical readouts.

\subsection{Auditability in Clinical Neuroscience}

The ability to audit intermediate transformations is relevant to scientific interpretation, even when no clinical translation claim is made.

For potential clinical use, a probability estimate alone is insufficient; the provenance, reliability, and limitations of the contributing measurements also matter. Post-hoc methods such as SHAP, attention visualization, and gradient saliency describe aspects of a learned model's function, while predefined temporal, dynamical, and topological variables describe the staged representation itself. These approaches answer different questions, and neither constitutes clinical validation on its own.

Post-hoc attributions and predefined intermediate summaries answer different questions. Neither reveals the latent physiological data-generating process automatically. A named variable is inspectable, but its fidelity, stability, physiological correspondence, and usefulness still require validation.

ARSPI-Net is designed around inspectable intermediate quantities: membrane trajectories and spike rasters; binned spike counts tied to fixed response intervals; per-channel embeddings; seven named dynamical descriptors; and graph-topological metrics. Some quantities have physical units, while others are dimensionless summaries or learned projections. Their value is that their construction is explicit and can be tested against conventional ERP measures; the dissertation does not evaluate clinician use.

\subsection{Spiking Neural Networks as a Temporal Representation}

Spiking neural networks (SNNs) represent activity through spike events generated by internal state dynamics \cite{maass_networks_1997}. In software, those events may be updated on a discrete simulation clock, as they are here. This representation creates two engineering opportunities.

The first is \emph{temporal expressivity}: spike timing, counts, and transient state evolution can be summarized explicitly instead of reducing the input immediately to a static amplitude vector. The dissertation tests whether that choice is useful on controlled signals and trial-averaged ERPs; it does not assume that spikes reproduce the biological encoding of EEG.

The second is \emph{potential deployment efficiency}: published hardware studies report low energy per event on particular neuromorphic chips \cite{davies2018loihi,schuman2022neuromorphic}. Per-operation values from different hardware and workloads are not directly comparable system-level energy measurements. The present NumPy implementation is not deployed on neuromorphic hardware, and no energy advantage is measured.

Among SNN models, the Liquid State Machine separates a fixed nonlinear temporal transformation from a downstream readout. Chapter~\ref{chap:dynamics} analyzes computational summaries of the driven reservoir. Those summaries are not biomarkers.

\subsection{The Graph Question: When Does Inter-Channel Structure Help?}

The LSM processes channels independently in this implementation. A graph layer is investigated as one way to model relationships among channel features. Because scalp channels are mixed sensor measurements, graph edges are modeling choices rather than direct functional organization.

Message-passing graph neural networks are used in multichannel EEG analysis \cite{Zhong2020EEGEmotionGCN,Song2018EEGEmotionDGCNN}. Their behavior depends on node features, graph construction, sample size, training, and the task. On the specific small, thresholded graphs tested here, smoothing may reduce feature contrast.

The dissertation compares several graph operators with a non-propagated representation and separately reports graph summaries. The results are limited to the tested cohort and implementation, and Chapter~\ref{chap:arspi_net} documents why the archived graph construction requires a fold-safe common-basis rerun.


\section{Research Objectives and Contributions}

This dissertation studies the \emph{Affective Reservoir-Spike Processing and Inference Network (ARSPI-Net)}, a staged pipeline comprising a fixed spiking reservoir, temporal spike-count coding, dimensionality reduction, conventional classifiers, and graph/dynamical summaries \cite{lane_arspi_2023,lane_arspi_2024}. The central engineering thesis is modest: exposing and testing each transformation can reveal where the pipeline is stable, useful to a decoder, or unsupported. The work does not recover latent brain dynamics, establish intrinsic clinical explanations, or demonstrate diagnostic utility.

\subsection{Research Objectives}

The specific research objectives are:
\begin{enumerate}
    \item To characterize a fixed LIF-like reservoir under controlled synthetic inputs and apply the specified transformation to SHAPE ERP matrices.
    \item To compare spike-train summaries and dimensionality-reduction choices for downstream decoding.
    \item To compare selected graph-propagation and non-propagation strategies while documenting the assumptions of graph construction.
    \item To define repeatable computational summaries of the driven reservoir response and evaluate their descriptive associations.
    \item To examine associations between reservoir-dynamical and graph-derived summaries, without treating them as physiological structure--function measurements.
\end{enumerate}

\subsection{Specific Contributions}

The dissertation makes the following specific contributions:
\begin{enumerate}
    \item \textbf{An auditable staged implementation.} The source defines the reservoir update, spike-count code, model summaries, and analysis sequence so that each computational transformation can be inspected.
    \item \textbf{Controlled representation experiments.} Synthetic experiments compare operating points and spike-train encodings. Results are tied to the tested generator and do not establish EEG generalization.
    \item \textbf{A critical real-data graph analysis.} The archived SHAPE runs compare graph and non-graph variants. The post-defense audit identifies global PCA, incompatible per-channel bases, and a full-epoch versus early-window BSC$_6$ mismatch. The corrected source specifies one fold-fitted common basis and the early window; affected numerical results remain descriptive pending rerun.
    \item \textbf{Reservoir-response summaries.} Seven computational summaries describe spike activity and temporal structure. Their exploratory condition effects are small, clinical-label estimates are near chance, and physiological meaning remains unvalidated.
    \item \textbf{Explicit inferential boundaries.} The dissertation distinguishes prospective from transductive panel estimates, nominal from multiplicity-controlled tests, internal resampling from replication, and model-derived associations from neural biomarkers.
\end{enumerate}

\subsection{Dissertation Organization}

The dissertation follows the pipeline from controlled reservoir experiments to the SHAPE analyses and their limitations.

Chapter~\ref{chap:background} reviews the signal-processing, dynamical-systems, neuromorphic, and graph-learning background, including the theoretical foundations of reservoir computing, graph signal processing, and the limits of message-passing on small dense graphs.

Chapter~\ref{chap:lsm} characterizes the reservoir on controlled synthetic inputs and selects an operating point carried into later chapters.

Chapter~\ref{chap:embeddings} compares spike-based encodings and dimensionality reduction on synthetic data. Its transfer to SHAPE requires the fold-local corrections described in Chapter~\ref{chap:arspi_net}.

Chapter~\ref{chap:arspi_net} applies the archived workflow to SHAPE ($N=211$), compares graph strategies, and documents the post-defense reproducibility audit. Global-PCA and full-epoch BSC$_6$ results are descriptive legacy estimates, and clinical graph screens do not support biomarkers.

Chapter~\ref{chap:dynamics} defines reservoir-response summaries and reports exploratory condition and clinical-label comparisons.

Chapter~\ref{chap:synthesis} examines descriptive correlations between temporal and graph-derived summaries and their limitations.

Chapter~\ref{chap:conclusion} consolidates contributions, outlines the boundaries of the present work, and proposes future directions including neuromorphic hardware deployment and single-trial EEG processing.

\subsection{Scope and Delimitations}

The dissertation is an electrical-engineering study of a staged representation pipeline for complex physiological time series. Its supported contributions concern implementation, controlled representation experiments, auditability, and explicit characterization of failure modes. It does not produce clinically meaningful characterizations of individual patients or validate neural-response mechanisms.

Several delimitations bound the scope. No clinician study, explanation-fidelity benchmark, prospective prediction study, independent cohort replication, source-localized EEG analysis, neuromorphic hardware deployment, or physical energy measurement is included. The SHAPE inputs are trial-averaged; electrode names are unavailable; clinical labels overlap; and the source package omits raw EEG. The attention-profile permutation test is negative, graph values inherit the archived PCA limitations, and reservoir summaries are model functionals rather than cortical measurements. These limits preclude clinical decision-support, biomarker, and deployment claims.
